\section{Related Work}
\label{sec:related}

\paragraph{Tool behavior and indirect attacks.}
ToolEmu evaluates risks in emulated environments, and R-Judge studies risk awareness in recorded interactions \citep{ruan2024toolemu,yuan2024rjudge}. Agent-SafetyBench covers varied environments and behavioral failures \citep{zhang2024agentsafetybench}. AgentHarm starts from malicious user requests and checks harmful task completion \citep{andriushchenko2025agentharm}. InjecAgent and AgentDojo instead place adversarial instructions in lower-trust inputs; AgentDojo checks task utility separately from attack success \citep{zhan2024injecagent,debenedetti2024agentdojo}. MT-AgentRisk studies how harmful tool use develops over multiple turns \citep{li2026unsafer}. SkillHarm and SkillSafetyBench extend evaluation to skills and related artifacts \citep{ning2026skillharm,jin2026skillsafetybench}. These works inform our attack construction and effect checks. We focus on how consultation and later actions interact under different human replies.

\paragraph{Human collaboration and authority.}
HiL-Bench measures task completion, blocker recall, and question precision when agents need missing information \citep{trinh2026hilbench}. HAS-Bench varies human participation and evaluates clarification, feedback use, control, and safety, including approval of protected actions \citep{wu2026hasbench}. HumanAgencyBench evaluates six aspects of agency support, such as clarification, correcting misinformation, and deferring important decisions \citep{sturgeon2025humanagencybench}. OverEager-Bench studies unauthorized scope expansion during benign coding tasks \citep{qu2026overeager}. Thus, human control and authorization are already evaluation targets. Our narrower question is how agents handle indirect attacks before and after feedback that differs in validity and scope. The planned factual-error and partial-approval tracks would extend the current four-reply core. Appendix~\ref{app:comparison} compares the evaluated settings without equating model counts with framework counts.

\paragraph{Defenses.}
ToolShield derives guidance from tool exploration, while Task Shield checks whether instructions and calls serve the user's task \citep{li2026unsafer,jia2024taskshield}. CaMeL separates trusted control from untrusted data \citep{debenedetti2025camel}. Progent represents privileges as symbolic policies and requires approval for policy expansion \citep{shi2025progent}. SafeAgent combines runtime control with recovery options, including replanning and human approval \citep{liu2026safeagent}. Task alignment, scoped privileges, and recovery therefore have substantial prior work. Our contribution connects a reply-sensitive benchmark to a test of original-authority preservation and task recovery. We compare controller versions and ablations within this benchmark, without claiming a head-to-head advantage over these defenses.
